\section{Prototype Design and Implementation}
\label{sec:prototype}

\subsection{Technology Stack}

Table~\ref{tab:stack} summarises the technology stack.  The system is
implemented entirely in Python 3.12 on both macOS and Linux (x86-64/ARM64), and
is fully compatible with Linux CUDA environments for GPU-accelerated
training.

\begin{table}[t]
\centering
\caption{Technology stack.}
\label{tab:stack}
\begin{tabular}{ll}
\toprule
Component & Library / Service \\
\midrule
Deep learning       & PyTorch 2.2.2 \\
Graph neural nets   & PyTorch Geometric 2.5.1 (SAGEConv, GCNConv, GATConv) \\
Language model      & Hugging Face Transformers 4.41.0 (CodeBERT) \\
Data integration    & DuckDB 0.10 \\
Tabular processing  & pandas, NumPy $<$2.0 \\
ML utilities        & scikit-learn 1.4 \\
Transaction source  & Google BigQuery (public Ethereum dataset) \\
Source-code source  & Etherscan API v2 \\
Fraud labels        & Forta Network labelled-datasets \\
Visualisation       & Matplotlib, seaborn \\
\bottomrule
\end{tabular}
\end{table}

\subsection{Dataset Characteristics}

After the full pipeline executes, the resulting graph contains
\num{85680} unique Ethereum addresses connected by \num{239742}
transactions.  Of these addresses, \num{304} carry both verified Solidity
source code and an explicit fraud/benign label---the eligible subset on
which all models are trained and evaluated.  Within this subset, 95
contracts are labelled as fraud (31.3\%) and 209 as benign.  The
70/15/15 stratified split produces 212 training, 45 validation, and 47
test samples.  Class weights are computed from the training split:
$w_0 = 1.0$ (benign), $w_1 \approx 2.37$ (fraud).

\subsection{Code Structure and Reproducibility}

The codebase is organised into nine top-level modules under \texttt{code/}:

\begin{enumerate}
    \item \texttt{config.py}: single source of truth for paths,
    hyperparameters, and API endpoints.
    \item \texttt{data/addresses/}: the two seed arrays
    (\texttt{BENIGN\_ADDRESSES}, \texttt{FRAUD\_ADDRESSES}).
    \item \texttt{src/collection/}: Etherscan v2 source-code fetcher,
    Forta label loader, BigQuery transaction reader.
    \item \texttt{src/preprocessing/}: address standardiser, DuckDB
    join engine, PyG graph constructor with eight-dimensional node
    features.
    \item \texttt{src/datasets/}: \texttt{ContractDataset} pairing
    Base64-decoded Solidity tokens with graph node indices.
    \item \texttt{src/models/}: \texttt{GraphEncoder} (GraphSAGE),
    frozen \texttt{load\_codebert}, and \texttt{FusionClassifier}.
    \item \texttt{src/training/}: AdamW + cosine-annealing loop
    with early stopping on validation fraud F1, gradient clipping,
    and best-checkpoint persistence.
    \item \texttt{src/evaluation/}: \texttt{evaluate\_model} returning
    precision, recall, F1, AUC-ROC, AUC-PR, and MCC; ablation
    classifiers for GraphSAGE-only, GCN-only, GAT-only, CodeBERT-only,
    attention-based fusion, and Random Forest baseline.
    \item \texttt{src/visualization/}: confusion-matrix and loss-curve
    plots for the paper figures.
\end{enumerate}

A \texttt{scripts/} directory contains eight numbered scripts (01--08),
one per pipeline step, each writing its outputs to a deterministic path
under \texttt{data/} so that later steps can be re-run independently.
The seed value 42 is set at both the Python and PyTorch levels before
any randomised operation, ensuring full reproducibility given the same
hardware and library versions.
